\documentclass[11pt]{article}
\usepackage[margin=.73in]{geometry}
\usepackage{lmodern,amsmath,amssymb,booktabs,graphicx,microtype,fancyhdr,xurl,hyperref}
\hypersetup{colorlinks=true,urlcolor=blue,linkcolor=black}
\pagestyle{fancy}\fancyhf{}\lhead{HPS GPR: understanding the 2021 signal extraction}\rhead{v6.3.6}\cfoot{\thepage}
\setlength{\headheight}{14pt}\setlength{\parindent}{0pt}\setlength{\parskip}{6pt}\setlength{\emergencystretch}{2em}
\newcommand{\Ah}{\widehat A}\newcommand{\sh}{\widehat\sigma}
\begin{document}
{\LARGE\bfseries Understanding signal extraction\par from the 2021 10\% sample}\par
{\large HPS Gaussian-process background analysis}\par
24 September 2026\hfill Version 6.3.6

\textbf{The result.} Moving the Gaussian signal template toward the reconstructed signal-MC core improves the fitted response. The Gaussian still represents only part of the asymmetric signal distribution. To make statements about the \emph{full selected signal yield}, the study therefore calibrates the fit statistic using simulated experiments with a known injected signal yield.

This note explains that procedure and the evidence supporting it. It is a presentation revision of v6.3.5: the main-body numerical ensembles, fitted results and observed scans are unchanged. Appendices A--C add v16 TC/UC core-shift studies; D--H add the v6.3.7 signal-template and window studies; I isolates signal leakage and its effect on the fit. The signal-MC catalogue and core-shift comparison from v6.1 are included so that the choice of Gaussian center can be understood here without consulting another note.

\textbf{Why the distinction matters.} There are three separate questions. Where is the reconstructed signal core? How much does the fitted Gaussian yield change when signal is added? Which injected yields are compatible with the observed fit result? A shifted Gaussian addresses the first question, an injection study measures the second, and toy-calibrated tests answer the third. Moving the Gaussian alone does not solve all three.

For signal-MC injections on the GP mean background source at $A=3s_0$ (where $s_0$ is the yield-error scale from background-only pilot fits), the added-signal response ranges from 0.208--0.422 for the central mass Gaussian and 0.265--0.833 for the shifted Gaussian across 60--240 MeV. A response of 0.5 means that adding 1,000 selected signal events increases the fitted Gaussian yield by about 500 events. It is not a detector efficiency. When the injected signal is itself a matched Gaussian, the shifted fit returns 0.963--0.980 of the added yield.

\textbf{What is usable now.} The finite-grid toy calibration provides a defined, testable procedure under the specified background sources and signal-MC shape. At $A=5s_0$, independent evaluation experiments retain the true injected yield in 81--96\% of cases for the GP mean source and 84--93\% for the functional form source. Each fraction uses 100 experiments. These are conditional validation results with finite statistical precision, rather than a guarantee of exactly 90\% coverage at every mass or for an unknown background shape.

\textbf{Reading guide.} Sections 1--2 explain the signal shape and fit. Sections 3--7 distinguish background-only bias, signal response and yield calibration. Section 8 explains the upper-limit construction. Sections 9--11 show the observed 2021 and combined results. Section 12 checks the cost of changing the fit window. Section 13 records the scope and reproducible sources.

The main search and calibration domain is 60--240 MeV. The 260 MeV signal-MC sample is shown only as a shape comparison. The report does not establish a global discovery probability or a detector-qualified physical exclusion.
\clearpage\section*{1. Why move the Gaussian center?}

A narrow generated resonance need not reconstruct as a symmetric Gaussian centered at its generated mass. Detector response and selection change the reconstructed distribution. In the supplied target-constrained (TC) signal Monte Carlo, the fitted core is displaced below the generated signal mass and a substantial tail extends toward higher reconstructed masses.

We compare two Gaussian extraction templates. The \textbf{central mass Gaussian} is centered at the tested signal mass $m$. The \textbf{shifted Gaussian} is centered at the empirical signal-core position $c(m)$:
\[
 c(m)=m-3.2243308693-2.2139928114\ln(m/150),\qquad 60\le m\le240\ \mathrm{MeV}.
\]
Here $m$ and $c(m)$ are in MeV; $m/150$ means $m/(150\,\mathrm{MeV})$. This is a logarithmic relation for the \emph{core displacement}, not a logarithmic probability model for the entire signal distribution.
\begin{center}\includegraphics[width=\linewidth,height=3.3in,keepaspectratio]{../figures/intro_core_shift_comparison.pdf}\end{center}
{\small \textbf{Figure 1.} \textbf{How to read the figure.} Left: the v6.1 TC signal-MC core displacement and the logarithmic relation used in this study. Left error bars are the standard deviation of the fitted center over 32 signal-MC bin-resampling replicas, often smaller than the markers; they are not the 100-toy mean errors used later. Right: the HPS Internal image supplied with this revision, reproduced unchanged. Its vertical coordinate is fitted Gaussian mean minus generated mass; negative values indicate a lower reconstructed center. The blue and red curves are the two sample categories named in that image. Its error definition and reconstruction/selection equivalence to the present input were not supplied, so the right panel provides qualitative context, not a combined fit. The logarithmic curve on the left is the one used here.}\par\medskip

Only the Gaussian center changes between the main fit choices. Both retain the same \textbf{reference analysis resolution},
\[
 \sigma_{\rm ref}=0.00184825-0.001375m+0.085875m^2,
 \quad m\ \hbox{and}\ \sigma_{\rm ref}\ \hbox{in GeV}.
\]
The signal-fit window and the region excluded from background training are both centered on that Gaussian and extend to $\pm2.25\sigma_{\rm ref}$. The fitted signal-MC core width used to compare shapes on the next pages is a different quantity.
\clearpage\section*{1. Signal-MC catalogue: examine each sample first}
\begin{center}\includegraphics[width=\linewidth,height=6.65in,keepaspectratio]{../figures/intro_signal_mc_catalogue.pdf}\end{center}
{\small \textbf{Figure 2.} \textbf{How to read the figure.} Each panel shows one generated signal-MC sample with the central mass Gaussian and the shifted Gaussian used for extraction. Both Gaussians have the reference analysis width, not the fitted signal-MC core width. The logarithmic vertical scale exposes tails over several orders of magnitude. These are signal-MC probability densities, not background toys; their full normalization is retained when only a local mass range is shown. The 260 MeV sample is a shape control outside the 60--240 MeV calibration domain and has no logarithmically shifted curve. No new fit was performed for this revision.}\par\medskip

The catalogue makes two features visible together: the central peaks move relative to their generated masses, and the selected signal is not fully Gaussian. An extraction template can locate the core well while still returning less than the full selected signal yield.
\clearpage\section*{1. Can one shape describe the signal cores?}
\begin{center}\includegraphics[width=\linewidth,height=4.6in,keepaspectratio]{../figures/intro_common_core_shape.pdf}\end{center}
{\small \textbf{Figure 3.} \textbf{How to read the figure.} The signal-MC distributions are aligned by their individual fitted core centers and widths. The horizontal coordinate is $u=(m_{\rm rec}-c_{\rm core})/\sigma_{\rm core}$, so a unit of $u$ is one fitted core width. The comparison tests whether the central shapes look similar after alignment; it does not imply that their full tails or acceptances are identical. Left: each curve is normalized within $|u|<2$ to compare core shapes. Right: the original core and tail probabilities are retained. The dashed black common shape is the equal-mass average of the ten aligned empirical distributions, not an analytic Gaussian fit. The gray dotted reference is a standard Gaussian normalized within $|u|<2$. These curves reuse the saved v6.1 histograms and core fits; no toy error bars are shown.}\par\medskip

\textbf{Answer.} A common empirical shape describes the cores reasonably well from about 80 MeV upward. Its peak is approximately Gaussian, with heavier shoulders than a standard Gaussian. The tails remain visibly different, and the low-mass sample needs separate treatment. Thus a common core description is a useful guide to the Gaussian center, but is insufficient to calibrate the total signal yield.

\textbf{The low-mass qualification.} The supplied independent signal-MC samples are generated at 60, 80, 100, \ldots, 260 MeV. There is no independently simulated 65 MeV sample in this catalogue. The fitted core of the 60 MeV sample is about 58.8 MeV. The acceptance-loss concern near 65 MeV is therefore illustrated here by the actual 60 MeV input, not by relabeling a sample. Its distorted distribution is consistent with the low-mass geometric acceptance turn-on discussed in the signal study; this plot does not isolate the cause. These selected distributions alone do not measure a separate acceptance value at 65 MeV.

The saved v6.1 diagnostic puts 32.39\% of its stored-range probability within two fitted core widths, compared with 60.73--71.47\% for 80--240 MeV. The figure also retains the tiny overflow probability (below 0.016\% in these samples). This helps explain why the 60 MeV response is particularly small in the injection study: fitting its central region with a Gaussian cannot account for the large fraction of selected signal outside that region. Subsequent sections measure that effect directly rather than assuming it away.
\clearpage\section*{2. What the Gaussian fit measures}
\begin{center}\includegraphics[width=\linewidth,height=4.35in,keepaspectratio]{../figures/shapes.pdf}\end{center}
{\small \textbf{Figure 4.} \textbf{How to read the figure.} Black: the selected signal-MC distribution used to inject events. Red: the central mass Gaussian. Blue: the shifted Gaussian, also called the \emph{log core Gaussian} because its center follows the logarithmic law in Section 1. Both Gaussians use the same reference analysis resolution. All curves retain their full probability normalization; the displayed mass range omits some tails without renormalizing what remains. These are the same extraction prescriptions as in the catalogue; the linear vertical scale emphasizes the core mismatch.}\par\medskip

Let $A$ be the expected number of \textbf{full selected signal events}. This counts the whole selected signal distribution, including events outside the analysis support and fit window. Let $\Ah$ be the signed Gaussian yield returned by the fit, and $\sh$ its reported standard error. Negative fitted yields are allowed so that background fluctuations are treated symmetrically; expected bin counts must remain positive.

The background is estimated with Gaussian-process regression (GPR) from the sidebands, meaning the bins outside the excluded signal region. Each toy or observed spectrum supplies its own sideband counts. The GP kernel settings remain fixed, but its predicted background mean and covariance are recomputed for that spectrum. The likelihood uses Poisson bin counts and correlated background nuisance parameters; the error $\sh$ comes from the curvature of the fitted profile likelihood.

Two effects reduce $\Ah$ relative to the full signal yield. First, the Gaussian fits only a projection of an asymmetric signal distribution. Second, signal in the training sidebands can be absorbed into the estimated background. With the background known exactly, the separate window study still finds response 0.378--0.857. The corresponding GP response divided by that control is 0.703--0.974. Moving the center helps, but neither effect is removed automatically.
\clearpage\section*{2. How the simulated experiments are organized}

A \textbf{toy experiment} is a simulated binned count spectrum. Background counts are drawn independently in each bin from Poisson distributions with fixed expected counts. We use two background sources:

\textbf{GP mean background source.} The expected bin counts are the pinned full-support GP arithmetic mean. We fluctuate counts around that fixed model; we do not draw a new GP background function for each toy.

\textbf{Functional form background source.} The expected bin counts come from the archived anchored \texttt{fSigPowExpQ} function. This alternative tests sensitivity to one specified change in background shape. Its earlier local assessment does not establish that it describes every part of the data globally.

Signal events are added according to the fixed selected signal-MC probabilities, including training sidebands and outside-support categories. At fixed expected signal yield $A$, their total fluctuates according to a Poisson distribution. These are \emph{not} injections with an exactly fixed event total. Neither the signal MC nor the Gaussian is renormalized inside the fit window.

\textbf{Three separate sets of experiments.} First, 100 GP-source background-only \emph{pilot} experiments define $s_0(m)$, the mean reported yield error from the shifted Gaussian. This fixes the signal-strength unit $A=z s_0$. The label $z$ therefore specifies injected yield in pilot-error units; it is not a discovery significance.

Next, each source has 100 new \emph{calibration} backgrounds. These determine the background-only bias, the response, and the reference distributions used by the yield tests. Those choices are frozen. Finally, each source has another 100 \emph{independent evaluation} backgrounds to test the frozen calibration. These evaluation experiments were called ``heldout'' in earlier versions: they were simply not used to determine the calibration.

Calibration uses $z=0,1,2,3,4,5,6,8,10,12,16,20,24$; evaluation uses $z=0,1,3,5$. Gaussian signal controls use $z=1,3,5$ on the GP mean source. The same expected yield is used for both Gaussian centers and both background sources.

Within a source and cohort, the same background experiment is reused across masses, strengths and fit choices. The two Gaussian fits also receive identical signal counts. Such comparisons are paired: their differences remove common fluctuations. Different sources and cohorts have independent streams. There are 500 main background experiments, with many fits to each, not 76,000 independent experiments.
\begin{center}\small\begin{tabular}{lrr}\toprule
Main calculation & Attempted & Valid\\\midrule
Pilot fits & 2,000 & 2,000\\
Calibration fits & 52,000 & 52,000\\
Evaluation fits / truth profiles & 22,000 / 22,000 & 22,000 / 22,000\\
Gaussian asymptotic limit endpoints & 16,000 & 16,000\\
\bottomrule\end{tabular}\end{center}\textbf{Uncertainty.} Means have sample-standard-error bars: the sample standard deviation divided by $\sqrt{100}$. Bootstrap calculations resample whole toy experiments 2,000 times, keeping masses, injected strengths and analysis choices together; the two background sources use separate streams. Widths have bootstrap uncertainty. Fractions retain all 100 trials and use exact two-sided 95\% Clopper--Pearson intervals. The caption specifies which of these bars appears here.\clearpage\section*{3. Measure background-only bias before adding signal}
\begin{center}\includegraphics[width=\linewidth,height=4.6in,keepaspectratio]{../figures/null.pdf}\end{center}
{\small \textbf{Figure 5.} \textbf{How to read the figure.} Both columns are background only ($A=0$). Left: calibration mean pull $\mu_0=\overline{\Ah_0/\sh_0}$. Right: mean of $\Ah_0/\sh_0-\mu_0$ in a separate evaluation cohort. Top: GP mean background source; bottom: functional form background source. Bars are sample standard errors of means from 100 experiments. The right-hand bars condition on the frozen left-hand calibration estimate. The uncertainty conventions are also printed on the figure.}\par\medskip

A background-only fit can return a nonzero average signal yield even though no signal was injected. We call this the \textbf{background-only fitted-yield bias}, $\delta=\overline{\Ah_0}$, measured in events. The \textbf{background-only mean pull}, $\mu_0=\overline{\Ah_0/\sh_0}$, instead measures displacement in units of the fit's own error. They are different averages and are not interchangeable.

Subtracting the frozen $\mu_0$ centers a pull diagnostic. Subtracting $\delta$ corrects an additive yield bias. Neither operation restores a missing fraction of injected signal. Nor does subtracting a mean change the width of a distribution.

For the shifted Gaussian, independent background-only mean pulls after centering range from $-0.174$ to $0.227$ for the GP mean source and $-0.187$ to $0.160$ for the functional form source. The corresponding uncentered pull widths are 0.832--1.131 and 0.893--1.121. Finite calibration and evaluation samples leave visible fluctuations; zero mean and unit width are targets to check, not assumptions to impose.
\clearpage\section*{4. Signal response before and after calibration}
\begin{center}\includegraphics[width=\linewidth,height=5.85in,keepaspectratio]{../figures/recovery_nominal.pdf}\end{center}
{\small \textbf{Figure 6.} \textbf{How to read the figure.} Signal-MC injections on the GP mean background source. Columns use $A=z s_0$ with $z=1,3,5$. Top: fitted-to-injected yield ratio $\overline{\Ah/A}$. Middle: change caused by adding signal, $\overline{(\Ah_{s+b}-\Ah_b)/A}$, with the same background experiment in both fits. Bottom: calibrated response $\overline{A_c/A}$, where $A_c=(\Ah-\delta)/R$. The independent calibration cohort supplies $\delta$ and the response $R$ measured at $z=3$. A value of one means that the estimator returns the injected full selected yield on average. Each point averages 100 independent evaluation experiments. Bars show sample standard errors, conditional on the frozen calibration table; uncertainty from estimating that table is assessed separately.}\par\medskip

The middle row isolates response to the added signal; the top row also contains background-only yield bias. The bottom row answers the requested calibration question directly. It applies the same frozen correction at all three strengths and tests it on new experiments. Its bars describe the evaluation mean for that frozen correction; they do not claim exact response under a different background source or signal shape.

A fixed subtraction of $\delta$ cancels in the middle-row difference, so it cannot repair a response below one. Dividing by $R$ addresses that scale difference. Section 7 explains the uncertainty of the resulting yield estimator; Section 8 uses full toy distributions for inference instead of relying on this simple correction.
\clearpage\section*{5. Repeat the test with a functional form background}
\begin{center}\includegraphics[width=\linewidth,height=5.85in,keepaspectratio]{../figures/recovery_functional.pdf}\end{center}
{\small \textbf{Figure 7.} \textbf{How to read the figure.} Signal-MC injections on the functional form background source. The rows and strengths have the same definitions as in Section 4: fitted-to-injected yield, signal-induced change divided by injected yield, and calibrated response $(\Ah-\delta)/(RA)$. Here the correction is derived from an independent calibration cohort of the same functional form source. Expected signal yields match the GP-source study, while source ensembles are independent. Each point averages 100 independent evaluation experiments. Bars show sample standard errors, conditional on the frozen calibration table; uncertainty from estimating that table is assessed separately.}\par\medskip

At $z=3$, the shifted-Gaussian signal-induced response is 0.265--0.834 for this source, close to the GP-source range. The fitted-to-injected ratio can differ more because the background-only biases differ. Agreement of the incremental response under two chosen sources is useful evidence, but does not test every possible background shape.

As a separate check, injected Gaussians centered on the logarithmic law give shifted-fit response 0.963--0.980 on the GP mean source. This is much closer to one than the signal-MC response. Signal-shape mismatch therefore matters even after the center is corrected. Gaussian controls have independent signal draws and receive no signal-MC response correction.
\clearpage\section*{6. What a pull says about bias and uncertainty}
\begin{center}\includegraphics[width=\linewidth,height=4.1in,keepaspectratio]{../figures/pulls.pdf}\end{center}
{\small \textbf{Figure 8.} \textbf{How to read the figure.} Independent signal-MC evaluation experiments at $A=5s_0$. The pull is $(\Ah-A)/\sh$: fitted yield minus the full selected injected yield, divided by the reported error. Top: its mean after subtracting the background-only mean pull $\mu_0$; bars are sample standard errors. Bottom: its sample standard deviation; bars are bootstrap standard errors from 2,000 whole-experiment resamples. Each point uses 100 experiments. The ideal reference lines are zero for the mean and one for the width.}\par\medskip

A negative mean pull says that the fitted Gaussian yield is below the injected full selected yield. A pull width near one only says that the fluctuation scale is roughly consistent with the reported error. It does not remove or excuse a displaced mean. This is why background-only centering can look satisfactory while the positive-signal test still shows a large bias.

The likelihood also provides a direct check. At the injected yield $A$, refit the background nuisance parameters and compute
\[
 q_A=2\{\ell(\Ah,\widehat\theta)-\ell(A,\widehat\theta_A)\},
\]
where $\ell$ is the log likelihood. The usual one-parameter Gaussian references accept $q_A\le1$ and $q_A\le3.841459$ for 68\% and 95\% likelihood sets. Here those sets are produced by a Gaussian-template model and checked against the full signal-MC yield. They need not contain that yield at the stated rates when the template response is biased.

The complete counts, always out of 100, and exact 95\% binomial intervals remain in \path{results/heldout_summary.csv}. That inherited filename denotes independent evaluation. The next two sections distinguish an approximate corrected-yield diagnostic from the toy-calibrated confidence construction.
\clearpage\section*{7. A calibrated yield and its error}

The corrected point estimate subtracts the background-only yield bias and divides by the measured response:
\[
 A_c=\frac{\Ah-\delta}{R},\qquad
 R=\overline{\frac{\Ah_{s+b}-\Ah_b}{3s_0}}\ \ \text{in the calibration cohort}.
\]
For example, if the fitted yield is 600 events, the calibrated background-only bias is 100 events and $R=0.5$, the corrected estimate is 1,000 selected events. Dividing by a small response amplifies the uncertainty as well as the yield.

Let $k_0={\rm SD}\{(\Ah_0-\delta)/\sh_0\}$ measure the background-only fluctuation scale. An approximate propagated variance is
\[
 \sigma_c^2\simeq\frac{k_0^2\sh^2+{\rm Var}(\delta)+A_c^2{\rm Var}(R)
             +2A_c\,{\rm Cov}(\delta,R)}{R^2}.
\]
The bootstrap estimates $\delta$ and $R$ together, retaining their covariance. The cross term has the displayed $+2A_c\,{\rm Cov}(\delta,R)$ form; its numerical contribution can have either sign.
\begin{center}\includegraphics[width=\linewidth,height=2.9in,keepaspectratio]{../figures/affine.pdf}\end{center}
{\small \textbf{Figure 9.} \textbf{How to read the figure.} Pull after yield calibration, $(A_c-A)/\sigma_c$, for independent signal-MC experiments at $A=5s_0$. A mean near zero indicates agreement between calibrated and injected yields for this source and strength. Bars are sample standard errors across 100 evaluation experiments, conditional on the frozen calibration table. Each toy's $\sigma_c$ includes the calibration-parameter terms in the equation above. The figure's bootstrap and sample-size definitions apply here.}\par\medskip

This first-order uncertainty estimate treats the response as locally constant and uses the background-only width scale. It omits uncertainty in the pilot scale, the chosen background source and the finite signal-MC template. It is a diagnostic of a corrected estimator, not a new likelihood or an exact confidence interval. The primary limits below therefore calibrate the fitted statistic at every tested injected yield; they are not obtained by simply dividing an existing limit by $R$.
\clearpage\section*{8. Turn the toy calibration into an upper-limit test}

The key question is: \emph{would an injected yield $A$ commonly give a fit result this small?} For each preselected yield on the calibration grid, compare the observed (or evaluation) fitted yield with 100 calibration fitted yields generated at that same true $A$:
\[
 p_A=\frac{1+\#\{\Ah_A^{\rm cal}\le\Ah^{\rm eval}\}}{101}.
\]
The numerator counts calibration results no larger than the result being tested, then adds one. Reject that yield when $p_A\le0.10$. A sufficiently small fitted yield is evidence against a large injected signal. This is a lower-side probability used for exclusion; Section 9 uses the opposite direction to ask about an excess above background.

Repeat this test at $z=0,1,2,3,4,5,6,8,10,12,16,20,24$, with $A=z s_0$. Retain the \emph{entire set} of accepted nodes. Its largest node is the displayed grid endpoint. We do not interpolate between simulated strengths. A missing interior node is recorded as a hole; acceptance of the largest node is recorded as right censoring, meaning that the grid does not reach the endpoint.

\textbf{An empty set is unresolved, not a zero physical limit.} If every node is rejected, the software stores endpoint zero together with an empty-set flag. If only zero is accepted, the grid resolves no behavior between zero and its first positive node. Neither case excludes all arbitrarily small positive signals. The accepted set and its flags must accompany the endpoint.

\textbf{Why this is called calibration.} The test uses distributions from signal-MC injections even though the fitted template is Gaussian. It therefore measures the distribution of the actual fitted statistic under the intended signal yield instead of assuming unit Gaussian-template response.

With continuous, exchangeable calibration and evaluation results, the add-one rank test rejects the true node in at most $10/101$ of experiments, averaged over both calibration and evaluation cohorts. A single frozen table can have a different acceptance probability. The order-statistic reference gives a central 95\% range of 83.60--95.10\% for that probability. This calibration-table variation is separate from the binomial error of checking it with 100 evaluation experiments; ties make the stated rank rule conservative.

For comparison, the study also records a toy-rank $CL_s$ diagnostic, $\min(1,p_A/p_{0,\mathrm{lower}})$, and the inherited Gaussian-template asymptotic $CL_s$ limit. The shared $1/101$ rank floor can limit the toy-rank diagnostic's sensitivity. The Gaussian asymptotic reference uses profile-likelihood approximations and does not automatically describe the full signal-MC yield.
\clearpage\section*{8. Check the procedure on independent experiments}
\begin{center}\includegraphics[width=\linewidth,height=4.45in,keepaspectratio]{../figures/limits.pdf}\end{center}
{\small \textbf{Figure 10.} \textbf{How to read the figure.} Shifted-Gaussian results, with calibration and evaluation using the same specified source. Upper panels: fraction of 100 independent experiments for which the injected $A=5s_0$ is accepted (or is below the continuous Gaussian asymptotic endpoint). Bars are exact two-sided 95\% Clopper--Pearson intervals. The dashed line marks 90\%. Lower panels: median upper endpoint in background-only experiments, divided by the pilot yield-error scale $s_0$; no uncertainty bars are shown on these medians. Toy-based endpoints are grid maxima, while Gaussian asymptotic endpoints are continuous.}\par\medskip

The toy-calibrated yield test retains the injected $A=5s_0$ in 81--96\% of GP-source evaluations and 84--93\% of functional-source evaluations. The corresponding Gaussian asymptotic limits contain the full signal-MC yield in only 4--72\% and 2--87\%. This failure is why improving the core center alone is insufficient for full-yield inference.

Every planned evaluation ID remains in the tables; no fits failed in the main run. The denominator is always 100. The variation among masses includes finite evaluation statistics and fluctuations of the frozen calibration tables, so the plot should not be summarized as exact 90\% coverage everywhere.

At zero injected yield, the stored numerical endpoint of an empty set also equals zero. Counting that endpoint as covering zero would hide the rejection of every tested node. The tables therefore retain both endpoint coverage and the separate, primary \emph{truth-acceptance} field. Use the latter to assess the grid test.
\clearpage\section*{9. What the observed 2021 spectrum shows}
\begin{center}\includegraphics[width=\linewidth,height=4.9in,keepaspectratio]{../figures/observed_2021.pdf}\end{center}
{\small \textbf{Figure 11.} \textbf{How to read the figure.} 2021 10\% observed data. Lines are Gaussian-template asymptotic results on a 1 MeV mass grid. Open circles are independent GP-source toy-calibrated results only at the simulated masses, 60--240 MeV in 20 MeV steps. Upper panel: yield limits converted to the inherited coupling display defined on the next page. Nonpositive grid endpoints are omitted from the logarithmic axis and retained with their flags in the tables. Lower panel: local excess probability. Lines use the profile-likelihood asymptotic approximation; open circles rank the observed fitted yield against background-only fitted yields. The line is not an interpolation of the points.}\par\medskip

For the toy-based local excess probability, count how many of the 100 background-only calibration experiments return a fitted yield at least as large as the observed one:
\[
 p_0=\frac{1+k}{101},\qquad k=\#\{\Ah_0^{\rm cal}\ge\Ah^{\rm obs}\}.
\]
This is the meaning of the former phrase ``upper null tail.'' The smallest resolvable rank is $1/101$, even when there are no exceedances. A binomial interval for $k/100$ estimates the underlying background-only tail probability; it is distinct from the add-one rank.

The smallest local asymptotic $p_0$ moves from 0.00249 at 78 MeV for the central mass Gaussian to 0.00346 at 80 MeV for the shifted Gaussian. These masses were selected by scanning the data. Their local probabilities do not include the look-elsewhere effect and are not global discovery probabilities.
\clearpage\section*{9. How fitted event yield is converted for display}

The fit fundamentally measures an event-yield parameter. The displayed coupling variable uses the inherited conversion
\[
 A=\epsilon^2 C_{2021}(m),\qquad
 C_{2021}(m)=\frac{3\pi m f_{\rm rad}}{2\alpha}\,\rho(m).
\]
Here $\alpha=1/137$, $f_{\rm rad}$ is the archived effective radiative fraction, and $\rho(m)$ is the saved observed event density averaged around the \textbf{tested signal mass} $m$. The density average uses $m\pm1.64\sigma_{\rm ref}$, including fractional overlaps with histogram bins. Mass and resolution are in GeV, and the density is events per GeV.

\textbf{The mass at which this conversion is evaluated does not move when the Gaussian center moves.} The central and shifted Gaussian fits at a given tested mass use the same $C_{2021}(m)$. This specifies precisely how the event density enters the conversion. It separates a change in extraction shape from a change in the yield-to-coupling normalization.

The implementation fits $\psi=\epsilon^2/10^{-8}$, so the fitted event yield is $\Ah=\widehat\psi\,C_{2021}(m)10^{-8}$. The raw coupling columns are electron-channel proxies. For consistency with earlier visible-channel plots, the displayed values above the muon-pair threshold are multiplied once by
\[
 1+\sqrt{1-4r}\,(1+2r),\qquad
 r=(105.6583745\,\mathrm{MeV}/m)^2,
 \quad m>211.316749\,\mathrm{MeV}.
\]
Below threshold the factor is one. This display convention does not independently validate branching fractions, efficiency or selection equivalence. It is therefore not a detector-qualified physical exclusion.

\textbf{How to interpret the two kinds of curve.} The continuous asymptotic line answers a question within the fitted Gaussian likelihood approximation. The open points use signal-MC toy distributions and the finite accepted-yield grid. Their agreement or disagreement is informative, but neither a sparse set of calibrated points nor a dense asymptotic scan establishes a globally calibrated discovery probability.
\begin{center}\small\begin{tabular}{lrr}\toprule
Gaussian choice & Dense minimum mass [MeV] & Local asymptotic $p_0$\\\midrule
Central mass Gaussian & 78 & 0.002488\\
Shifted Gaussian & 80 & 0.003464\\
\bottomrule\end{tabular}\end{center}\clearpage\section*{10. Combine campaigns while changing only 2021}
\begin{center}\includegraphics[width=\linewidth,height=4.7in,keepaspectratio]{../figures/observed_combined.pdf}\end{center}
{\small \textbf{Figure 12.} \textbf{How to read the figure.} Combined observed result with a common coupling parameter. The 2015 and 2016 models are unchanged; only the 2021 Gaussian center and its matched fit/training exclusion change. Lines are dense asymptotic references; open points are direct joint toy calibrations at simulated masses. All three campaigns contribute through 100 MeV, 2016+2021 through 180 MeV, and only 2021 above 180 MeV; dotted lines mark those boundaries. Empty and zero-only grid sets are omitted from the logarithmic endpoint panel. The coupling display is defined in Section 9.}\par\medskip

The joint calibration simulates the full common-coupling experiment, rather than combining separately corrected amplitudes. It uses 100 pilot experiments to set the coupling scale, 100 independent GP-source calibration experiments at the thirteen strengths, and another 100 independent evaluation experiments at $z=0,1,3,5$. The 2021 signal is drawn from signal MC; older campaigns retain their inherited signal models.

For the combined \emph{excess} probability, the ranked statistic is the signed profile-likelihood root. It is the square root of twice the best-fit log-likelihood improvement over background only, with the sign of the fitted coupling. Thus the background-only exceedance count uses likelihood evidence for an excess, not the standalone 2021 fitted yield. Exclusion still orders the fitted yield parameter.

The smallest combined local asymptotic $p_0$ changes from 0.00289 at 66 MeV to 0.000933 at 67 MeV. These mass-selected minima remain descriptive local references. Changing only the 2021 model does not independently recalibrate the older campaigns.
\clearpage\section*{11. Check the combined toy calibration}
\begin{center}\includegraphics[width=\linewidth,height=3.05in,keepaspectratio]{../figures/joint_validation.pdf}\end{center}
{\small \textbf{Figure 13.} \textbf{How to read the figure.} Fraction of independent joint evaluation experiments retaining the injected coupling at its grid node. Left: background only; right: signal plus background at $z=5$ in the joint pilot-defined coupling unit. Each point uses 100 experiments and exact two-sided 95\% Clopper--Pearson intervals. The two Gaussian choices share each whole experiment, preserving their correlation. Calibration and evaluation use the GP mean background source; campaign support follows Section 10.}\par\medskip

Across all joint evaluation cells, acceptance ranges from 78\% to 100\%; 260 of 8,000 accepted-node sets are empty. This reflects a finite, frozen calibration with discrete strengths. It does not establish exactly 90\% coverage for each mass and strength.

The observed joint inversions have five empty sets and five additional zero-only sets among twenty mass/fit-choice pairs. None accepts the largest grid node or has an internal hole. Their zero endpoints are unresolved continuous limits, not exclusions of arbitrarily small signals.
\begin{center}\small\begin{tabular}{lrrrr}\toprule
Mass [MeV] & Central $p_0$ & Shifted $p_0$ & Central grid result & Shifted grid result\\\midrule
60 & 0.9604 & 0.9208 & empty & empty\\
80 & 0.1584 & 0.0396 & 5.26e-06 & 5.26e-06\\
100 & 0.8416 & 0.9505 & zero only & empty\\
120 & 0.1980 & 0.3960 & 6.62e-06 & 1.32e-06\\
140 & 0.8515 & 0.9307 & zero only & empty\\
160 & 0.0594 & 0.1089 & 1.14e-05 & 5.69e-06\\
180 & 0.1683 & 0.8515 & 1.20e-05 & zero only\\
200 & 0.3663 & 0.6337 & 9.37e-06 & zero only\\
220 & 0.9505 & 0.8020 & empty & zero only\\
240 & 0.3762 & 0.3960 & 3.74e-05 & 2.81e-05\\
\bottomrule\end{tabular}\end{center}
Positive grid entries are endpoints in the coupling display of Section 9. The smallest joint empirical rank is $4/101=0.03960$ at 80 MeV for the shifted Gaussian, corresponding to three background-only exceedances. This finite sample does not empirically resolve the much smaller dense asymptotic minimum.
\clearpage\section*{12. Would a different fit or training window help?}

A wider excluded training region removes more signal contamination from the background estimate, but also removes background information. The relevant test is the uncertainty on the full signal yield after accounting for response, not signal containment alone.

The baseline fit and excluded training region both extend $2.25\sigma_{\rm ref}$ on each side of the shifted center. Three alternatives change both together to left/right half-widths $(2,2)$, $(2.5,2)$ and $(2,2.5)$ in units of $\sigma_{\rm ref}$. A fourth keeps the baseline fit window and widens only the excluded training region to contain the central equal-tail 95\% of selected signal MC.
\begin{center}\includegraphics[width=\linewidth,height=3.5in,keepaspectratio]{../figures/windows.pdf}\end{center}
{\small \textbf{Figure 14.} \textbf{How to read the figure.} Separate GP-source side study with 100 pilot and 100 independent evaluation backgrounds. Signal-MC injections use a common $A=3s_0$ and the same experiments across window choices. The plotted ratio compares $\mathrm{SD}(\Ah_0)/R$ for each choice with the baseline: a lower value means less full-yield noise. Bars are 95\% percentile intervals from 2,000 whole-toy bootstrap resamples, preserving mass and window-choice correlations; the pilot scale is fixed. These are sensitivity diagnostics, not calibrated upper limits.}\par\medskip

The three matched-window alternatives give ratios 0.983--1.014, mostly with intervals spanning one. The wider-left 100 MeV case is 0.983 [0.972, 0.997], one exploratory point among thirty comparisons. This does not establish a reliable improvement. The 95\% training exclusion worsens the noise by factors 1.15--3.54, despite reducing training leakage to about 5\%.

At 60 MeV that equal-tail exclusion spans roughly 51.0--212.25 MeV, leaving a large gap for background prediction. Keeping three external bins on either side checks geometry, not the reliability of that extrapolation. The baseline $\pm2.25\sigma_{\rm ref}$ prescription is retained.

The shape coordinate $u=(m_{\rm rec}-c)/\sigma_{\rm core}$ from Section 1 must not be confused with these widths in $\sigma_{\rm ref}$. Translating a boundary between the two requires the ratio $\sigma_{\rm core}/\sigma_{\rm ref}$.
\clearpage\section*{13. What is established, and how to reproduce it}

\textbf{Established within this study.} A shifted Gaussian improves extraction of the supplied signal MC. Background-only bias and incomplete signal response are different effects and can be measured separately. A frozen toy-calibrated yield test gives a defined conditional inference procedure that can be checked on independent experiments. The chosen baseline window remains appropriate among the tested alternatives.

\textbf{Scope of that conclusion.} The background sources, empirical signal-MC histogram and archived GP kernel settings are fixed. The selected TC signal-MC metadata identifies v13 inputs; it does not establish full v13/v16 selection equivalence or daughter association. Uncertainty in the finite signal-MC template, other possible background shapes, model selection and a global look-elsewhere correction are outside this ensemble. Earlier 2016 support/optimizer qualifications remain unresolved by these 2021 toys. A residual in one observed spectrum is not itself an ensemble bias estimate.

\textbf{Numerical evidence.} All main planned fits, truth profiles and Gaussian asymptotic endpoints passed the archived finite/positive, score and Hessian checks. The original validation includes seed replay, probability normalization, paired draws, frozen-calibration hashes and representative direct calculations. The main-body revision reruns no inference fits. The appendices add new v16 signal-MC core fits. Hash comparisons verify that the numerical inputs, cohorts, calibration and result tables are unchanged from the supplied v6.3.5 archive.

\textbf{Reproduction.} The package contains \texttt{source/report.tex}, the plot and narrative builders in \texttt{scripts/}, figure source data, and the original numerical checkpoints. The README gives the report-only rebuild commands. \texttt{protocol.json} deliberately retains the v6.3.5 scientific protocol and seed so that editorial versioning cannot be mistaken for a new experiment. \texttt{revision.json} describes v6.3.6. The source image is preserved unchanged with a SHA-256 digest. New and inherited QA records are identified separately.

The source keys \texttt{nominal}, \texttt{functional}, \texttt{pole} and \texttt{logshift} are preserved in machine-readable files for compatibility. In this note they mean GP mean background source, functional form background source, central mass Gaussian and shifted Gaussian, respectively. Inherited ``heldout'' filenames contain independent evaluation results. No terminology change alters the estimator or data.

\textbf{Related studies.} The v6.1 signal-MC study supplies the catalogue, core fits and logarithmic center law. The v6.2 and v6.3.1 matched-template studies generated and fitted signal MC; their closure does not by itself validate a Gaussian fit to signal-MC injections. The v6.3.2 study concerned controlled 2016 exposure changes. The main-body inference results are the frozen v6.3.5 results. The appendices provide the new v16 TC and UC signal-shape studies.

\textbf{Statistical references retained from v6.3.5.} The PDG statistics review discusses bias, confidence constructions and $CL_s$: \url{https://pdg.lbl.gov/2026/reviews/rpp2026-rev-statistics.pdf}. The profile-likelihood asymptotic reference is Cowan et al., \url{https://arxiv.org/abs/1007.1727}. The finite-rank rule actually used here is given explicitly in Section 8.
\clearpage\section*{Appendix A. Locate and verify the v16 signal samples}

The nearby v16 production contains both requested signal samples. The target-constrained (TC) and unconstrained (UC) categories are established by the stored vertex-type branch, not inferred from their directory names. Both are selected 2021 signal Monte Carlo; this appendix studies their reconstructed mass shapes without an observed-data fit.

The common directory is \path{/sdf/data/hps/physics2021/preselection/v16/}. Its \path{ap_signal_prompt_smeared} child contains TC vertices, while \path{ap_signal_prompt_dis_smeared} contains UC vertices. We read the latest \texttt{preselection} tree once per file, using \path{vertex./vertex.invM_} in GeV. No extra cuts, smearing or recentering are applied to the stored masses.
\begin{center}\small\begin{tabular}{lrr}\toprule
Verified quantity & TC & UC\\\midrule
Generated samples & 60--260 MeV, step 20 & 60--260 MeV, step 20\\
Source ROOT files & 24 & 24\\
Selected candidate rows & 17,775,543 & 21,089,537\\
Nonunit event weights & 0 & 0\\
Underflow / overflow & 0 / 1,205 & 0 / 1,481\\
Stored vertex type & TargetConstrained & Unconstrained\\
\bottomrule\end{tabular}\end{center}
Histograms use 0.1 MeV bins on 0--400 MeV. Their bin sums plus flow counts reproduce all selected rows. Generated truth masses agree with the sample labels within 0.003914 MeV. All 22 samples have finite reconstructed masses and positive unit weights. As in the main note, 260 MeV is a shape control; the shift relations and common-core comparisons use 60--240 MeV. There is no independent 65 MeV sample.

\textbf{These are different selected productions.} TC cutflow labels list electron/positron two-dimensional-hit thresholds of 8/10 and momentum sum above 2.8 GeV; UC labels list 6/8 and momentum sum above zero. Both display a vertex-$\chi^2<30$ step, but the TC count does not decrease there in any file, whereas the UC count does. The timing labels are 5.7 and 9.2 ns in both productions, so a v16 directory name alone does not certify equivalence to the observed-data selection.

The stored momentum-smear ratios equal one in TC and are generally nonunit in UC. These auxiliary branches do not by themselves establish the transformations already included in the stored mass. We therefore histogram that mass exactly once. A TC--UC shift difference can reflect vertex treatment, selection and smearing conventions together; it is not an isolated measurement of the effect of the vertex constraint.

\textbf{Reproducible extraction.} A bounded read-only S3DF pass streamed the mass and audit branches with one decompression worker. The saved provenance contains all file paths, sizes, modification times, tree cycles, cutflows and SHA-256 hashes of streamed branch payloads. These hashes identify the analyzed payloads, not the complete ROOT files. All core fits and plots were then made locally. The main inference results and their frozen inputs remain unchanged; no UC upper limit is calculated because a matching observed UC source histogram is not available.
\clearpage\section*{Appendix A. Use the same core definition for both samples}

The goal is to measure the reconstructed signal core without forcing its asymmetric tail into a single Gaussian. The same algorithm used for v6.1 is applied separately to every v16 TC and UC histogram.

First smooth the histogram only to locate its peak within three reference widths of the generated mass. Then fit the original, unsmoothed bin counts within $1.5\sigma_{\rm ref}$ of that peak with a bin-integrated Gaussian plus a nonnegative linear pedestal. The pedestal describes the local signal tail under the core; it is not an experimental background. The Gaussian center $c_{\rm core}$ and width $\sigma_{\rm core}$ are free. The mass displacement is $c_{\rm core}-m$.

The TC reference-resolution formula from Section 1 supplies the search and fit-span scale for both categories. In UC it is only a common numerical scale used to define the fit region, not an assumed or calibrated UC detector resolution. The UC Gaussian width is fitted freely. As a stability check, repeat the fit with half-spans $1.25\sigma_{\rm ref}$ and $2\sigma_{\rm ref}$.

\textbf{Uncertainty on a core position.} Each mass histogram is independently resampled 32 times by drawing a Poisson count in every bin with mean equal to that bin's observed signal-MC count. The error bar is the sample standard deviation of the 32 refitted centers. All 22 central fits, all 66 fit-span checks and all 704 resampling fits passed the optimizer, boundary and peak-location checks. These bars describe finite signal-MC statistics conditional on the selected production and core definition. They are not standard errors of 100 background toys, and they do not include selection, smearing or fit-definition uncertainty.

\textbf{Which shift relation describes the mass dependence?} Fit four simple functions by unweighted least squares to the ten core displacements from 60 to 240 MeV:
\[
 a,\qquad a+bx,\qquad a+b\ln[m/(150\,\mathrm{MeV})],\qquad a+bx+dx^2,
 \quad x=\frac{m-150\,\mathrm{MeV}}{100\,\mathrm{MeV}}.
\]
To compare them, omit one mass, fit the remaining nine, and predict the omitted displacement. Repeat for all ten masses and report the root-mean-square (RMS) prediction error in MeV. The preferred description has the smallest error among these four choices. This is a model-selection diagnostic, not independent validation or a guarantee outside the sampled domain. All coefficients describe $c-m$ in MeV.

\textbf{Can one shape describe the cores?} Align each histogram using its own fitted $c_{\rm core}$ and $\sigma_{\rm core}$. A core-only comparison normalizes each distribution within $|u|<2$, where $u=(m_{\rm rec}-c_{\rm core})/\sigma_{\rm core}$. A second panel retains full selected normalization to expose the tails. The common empirical curve gives equal weight to each generated mass; it is not a Gaussian fit. The quantitative comparison is the largest difference of normalized cumulative probabilities inside the core, with the tested mass omitted from the common curve.
\clearpage\section*{Appendix B. The v16 target-constrained (TC) core shift}
The v16 TC sample reproduces the earlier logarithmic core-shift pattern. The logarithmic relation gives the smallest omitted-mass prediction error of the four tested forms:\[c_{\rm TC}(m)-m=\left[-3.224325-2.213983\ln[m/(150\,\mathrm{MeV})]\right]\mathrm{MeV}.\]\begin{center}\includegraphics[width=\linewidth,height=4.25in,keepaspectratio]{../figures/../appendix/figures/tc_core_shift.pdf}\end{center}
{\small \textbf{Figure 15.} Upper panel: v16 TC fitted core displacement and four fitted descriptions; the original v13 TC relation supplies context. Lower panel: predicted minus measured displacement when each mass is omitted from its relation fit. Points use generated masses 60--240 MeV. Error bars on fitted centers are sample standard deviations from 32 independent Poisson-bin resamples per histogram; the prediction-error curves have no uncertainty bars. These descriptive shape comparisons do not establish sample independence or full selection equivalence. No significance is assigned to differences between productions.}\par\medskip
\begin{center}\small\begin{tabular}{lr}\toprule
Shift description & Omitted-mass RMS [MeV]\\\midrule
Constant & 1.0647\\
Linear & 0.2153\\
Logarithmic & 0.0652\\
Quadratic & 0.1399\\
\bottomrule\end{tabular}\end{center}Its omitted-mass RMS error is 0.0652 MeV. At the directly simulated masses, the largest v16--v13 TC core-center difference is only 0.000078 MeV. The histograms are very close but not identical. This comparison does not establish independent samples or full selection equivalence. The core-definition span check moves a center by up to 0.162 MeV, a separate effect from the statistical bars.\clearpage\section*{Appendix B. The TC signal-MC catalogue}
\begin{center}\includegraphics[width=\linewidth,height=6.3in,keepaspectratio]{../figures/../appendix/figures/tc_catalogue.pdf}\end{center}
{\small \textbf{Figure 16.} All eleven v16 TC samples. Black: selected signal-MC density, normalized to all selected candidates including overflow. Blue: local Gaussian plus nonnegative linear pedestal, drawn only across its fitted interval. Red dashed: the Gaussian component of that local fit, with its fitted amplitude and free width. These curves summarize the core; they are not the unit-area extraction Gaussians used in the main inference study. The logarithmic vertical scale shows the tails. The 260 MeV panel is shape-only and is excluded from the fitted mass-shift relations. No additional smearing or mass translation was applied.}\par\medskip
A good description of the local peak does not mean that the fitted Gaussian contains the full selected signal probability. The pedestal is a local fitting device; it is not extrapolated as a physical full-range signal model. The next page compares aligned cores while retaining the original tails in a separate panel.\clearpage\section*{Appendix B. How similar are the aligned TC cores?}
\begin{center}\includegraphics[width=\linewidth,height=3.8in,keepaspectratio]{../figures/../appendix/figures/tc_common_core.pdf}\end{center}
{\small \textbf{Figure 17.} v16 TC distributions aligned by each sample's fitted core center and width. Left: each curve is normalized within $|u|<2$; right: full selected normalization retains the original tail probabilities. Black dashed: equal-mass empirical common shape. Gray dotted: standard Gaussian conditioned on $|u|<2$. Individual curves are signal-MC histograms; no error bars are shown. The core-only normalization is used solely for this shape comparison, never to renormalize an injected signal yield. The comparison uses 60--240 MeV.}\par\medskip

From 80 to 240 MeV, the largest omitted-mass cumulative-probability difference inside the core ranges from 0.0067 to 0.0137. This supports a common empirical core description at the percent level in this specific diagnostic. It does not establish a universal full signal shape, common acceptance or an exact Gaussian core.

At 60 MeV the corresponding difference is 0.0450, and only 32.39\% of the full selected probability lies within two fitted core widths. The 80--240 MeV core fractions range from 60.73\% to 71.46\%. Thus the low-mass sample still needs separate treatment even after its center and width are aligned.
The near identity of the v13 and v16 TC core law supports the stability of the center prescription in the main note. It does not replace the original injection calibration, assess all selection changes, or propagate finite signal-template uncertainty into those limits. The existing TC inference results are retained unchanged. This appendix evaluates the shift and shape evidence rather than rerunning observed limits or the full injection calibration.
\textbf{Files.} The extracted histograms and per-file audit metadata are in \path{appendix/inputs/}; core fits, all resampling fits, alternative spans, prediction checks and shape metrics are in \path{appendix/results/}. The analysis and figure source is \path{appendix/scripts/analyze_shifts.py}. These data are sufficient to rebuild the appendix locally without another S3DF read.
\clearpage\section*{Appendix C. The v16 unconstrained (UC) core shift}
UC also reconstructs below the generated mass, but its displacement is smaller and has a different mass dependence. Among the four tested forms, the quadratic gives the smallest omitted-mass prediction error:\[c_{\rm UC}(m)-m=\left[-0.974838-0.791808x-0.296125x^2\right]\mathrm{MeV},\quad x=\frac{m-150\,\mathrm{MeV}}{100\,\mathrm{MeV}}.\]\begin{center}\includegraphics[width=\linewidth,height=4.25in,keepaspectratio]{../figures/../appendix/figures/uc_core_shift.pdf}\end{center}
{\small \textbf{Figure 18.} Upper panel: v16 UC fitted core displacement and four fitted descriptions; the original v13 TC relation supplies context. Lower panel: predicted minus measured displacement when each mass is omitted from its relation fit. Points use generated masses 60--240 MeV. Error bars on fitted centers are sample standard deviations from 32 independent Poisson-bin resamples per histogram; the prediction-error curves have no uncertainty bars. These descriptive shape comparisons do not establish sample independence or full selection equivalence. No significance is assigned to differences between productions.}\par\medskip
\begin{center}\small\begin{tabular}{lr}\toprule
Shift description & Omitted-mass RMS [MeV]\\\midrule
Constant & 0.5152\\
Linear & 0.1271\\
Logarithmic & 0.2482\\
Quadratic & 0.0422\\
\bottomrule\end{tabular}\end{center}The UC displacement runs from $-0.519$ MeV at 60 MeV to $-1.956$ MeV at 240 MeV. The quadratic prediction RMS is 0.0422 MeV, compared with 0.2482 MeV for a refitted logarithmic form. Thus the TC logarithmic correction should not be copied into UC. Changing the core-fit span moves a center by up to 0.088 MeV in this domain; the statistical bars do not include that variation.\clearpage\section*{Appendix C. The UC signal-MC catalogue}
\begin{center}\includegraphics[width=\linewidth,height=6.3in,keepaspectratio]{../figures/../appendix/figures/uc_catalogue.pdf}\end{center}
{\small \textbf{Figure 19.} All eleven v16 UC samples. Black: selected signal-MC density, normalized to all selected candidates including overflow. Blue: local Gaussian plus nonnegative linear pedestal, drawn only across its fitted interval. Red dashed: the Gaussian component of that local fit, with its fitted amplitude and free width. These curves summarize the core; they are not the unit-area extraction Gaussians used in the main inference study. The logarithmic vertical scale shows the tails. The 260 MeV panel is shape-only and is excluded from the fitted mass-shift relations. No additional smearing or mass translation was applied.}\par\medskip
A good description of the local peak does not mean that the fitted Gaussian contains the full selected signal probability. The pedestal is a local fitting device; it is not extrapolated as a physical full-range signal model. The next page compares aligned cores while retaining the original tails in a separate panel.\clearpage\section*{Appendix C. How similar are the aligned UC cores?}
\begin{center}\includegraphics[width=\linewidth,height=3.8in,keepaspectratio]{../figures/../appendix/figures/uc_common_core.pdf}\end{center}
{\small \textbf{Figure 20.} v16 UC distributions aligned by each sample's fitted core center and width. Left: each curve is normalized within $|u|<2$; right: full selected normalization retains the original tail probabilities. Black dashed: equal-mass empirical common shape. Gray dotted: standard Gaussian conditioned on $|u|<2$. Individual curves are signal-MC histograms; no error bars are shown. The core-only normalization is used solely for this shape comparison, never to renormalize an injected signal yield. The comparison uses 60--240 MeV.}\par\medskip

From 80 to 240 MeV, the largest omitted-mass cumulative-probability difference inside the core ranges from 0.0026 to 0.0079. This supports a common empirical core description at the percent level in this specific diagnostic. It does not establish a universal full signal shape, common acceptance or an exact Gaussian core.

At 60 MeV the corresponding difference is 0.0180, and only 34.48\% of the full selected probability lies within two fitted core widths. The 80--240 MeV core fractions range from 53.63\% to 68.74\%. Thus the low-mass sample still needs separate treatment even after its center and width are aligned.
The UC core is broader than the TC core in these samples, and its smaller shift has a different mass dependence. These conclusions concern the complete selected productions, whose selection and smearing information differ. The agreement in sign with the HPS Internal comparison in Section 1 is qualitative; that supplied image does not provide enough provenance for a quantitative combined fit.

No UC upper limit, local excess probability or signal-recovery calibration is computed. Those would require the matching observed UC histogram and an appropriate background and yield model. The appendix supplies the signal-shape evidence needed to design that later study.
\textbf{Files.} The extracted histograms and per-file audit metadata are in \path{appendix/inputs/}; core fits, all resampling fits, alternative spans, prediction checks and shape metrics are in \path{appendix/results/}. The analysis and figure source is \path{appendix/scripts/analyze_shifts.py}. These data are sufficient to rebuild the appendix locally without another S3DF read.
\clearpage\section*{Appendix D. v6.3.7: fit the signal-MC shape and test the window}

A Gaussian describes the narrow signal core, but it does not describe all selected signal events. This study asks whether using the measured core and tails improves the extraction of a known injected signal. It also asks how much mass range to fit and how much to remove from Gaussian-process (GP) background training. Better shape agreement alone does not guarantee a better upper limit: removing more sideband bins weakens the background prediction.

The inputs are verified v16 target-constrained (TC) signal Monte Carlo: stored reconstructed masses binned at 0.1 MeV, normalized to all selected candidates including overflow. The new templates use simulated masses from 80 to 240 MeV. The acceptance-turn-on sample at 60 MeV is not used to predict shapes in this interval. No extrapolation is made outside it. UC inference is not attempted because there is no matching UC background source histogram.

\textbf{Two different intervals.} Let $F$ denote the bins included in the signal likelihood, and $G$ the bins excluded from GP training. The coordinate is
\[
 u=\frac{m_{\rm rec}-c(m)}{\sigma_{\rm core}(m)}.
\]
Here $c(m)$ and $\sigma_{\rm core}(m)$ are predicted signal-core quantities. They are different from the reference analysis resolution used by the original shifted Gaussian. The proposed starting exclusion is $G_0=[-4,3]$ in $u$. To avoid using a count both to train the background and to fit the signal, this implementation requires
\[
 F\subseteq G,\qquad G=G_0\cup F\quad\hbox{when extending the starter fit.}
\]
Thus extending the fitted template by two core widths on each side changes both the fitted interval and the actual GP exclusion to $[-6,5]$. A guard-only extension, which removes training bins without fitting them, tests a different tradeoff. The GP kernel parameters remain frozen; its background mean and covariance are recomputed from each toy's remaining sidebands.

\textbf{What a fitted yield means.} The expected yield $A$ counts all selected signal events. Every injected toy uses the full signal-MC distribution, including tails in GP training bins and events outside the background spectrum. For a template probability $p_i$, the signal expectation in fitted bin $i$ is $A p_i$. If the fit retains probability $f_F$, its expected signal count is $A f_F$. We never renormalize the retained bins to unit probability. Shortening the fit therefore cannot make physical leakage disappear.

The GP mean background source generates Poisson count fluctuations about the archived mean spectrum. It does not draw a random GP function for each toy. These are conditional sensitivity studies using the selected signal production and the archived background model, not a new observed-data search. The main-body v6.3.6 observed fits and limits are unchanged.
\clearpage\section*{Appendix E. Predict an intermediate mass from its neighbors}

For a tested mass $m$ between simulated masses $m_L$ and $m_R$, set $t=(m-m_L)/(m_R-m_L)$. Interpolate the core center and width,
\[
 c_m=(1-t)c_L+t c_R,\qquad s_m=(1-t)s_L+t s_R.
\]
For each anchor, define its complete aligned cumulative distribution $C_j(u)=P_j[(m_{\rm rec}-c_j)/s_j\le u]$. The prediction is
\[
 F_m(x)=(1-t)C_L\!\left(\frac{x-c_m}{s_m}\right)
        +t C_R\!\left(\frac{x-c_m}{s_m}\right).
\]
Bin probabilities are differences of this cumulative distribution at bin edges. The weights are nonnegative and sum to one, so the construction preserves probability and moves both the core and the tails. The recorded overflow remains an unresolved upper-tail category; it is not assigned invented mass positions.
\begin{center}\includegraphics[width=\linewidth,height=3.65in,keepaspectratio]{../figures/../appendix_v637/figures/v637_neighbor_interpolation.pdf}\end{center}
{\small \textbf{Figure 21.} \textbf{How to read the figure.} The 90 MeV model interpolates the 80 and 100 MeV signal-MC samples with equal weights after aligning each fitted core. Both the fitted core location and width interpolate linearly; the complete aligned CDF carries the tails. The left panel resolves the core and the right panel resolves small tail probabilities on a logarithmic scale. The shaded interval is the proposed training exclusion [-4,3]. The 90 MeV curve is a conditional prediction, not closure against an independent sample. No uncertainty bands are shown.}\par\medskip

At 90 MeV, the fitted core is predicted at 87.895 MeV with width 1.745 MeV. Both aligned signal distributions receive weight one half. This implements the requested intermediate core and tails. It is related to the alignment-and-mixing idea in template morphing [1]; the core-fit parameters used here are not the full moments of that method.

The \textbf{common-shape comparator} instead averages the full aligned distributions at all retained anchor masses with equal weight. Its target center and width still come from the nearest two masses. Comparing it with neighboring-template interpolation isolates the benefit of local tail information.
\clearpage\section*{Appendix E. Predict samples that were removed from the model}
\begin{center}\includegraphics[width=\linewidth,height=5.0in,keepaspectratio]{../figures/../appendix_v637/figures/v637_omitted_mass_shapes.pdf}\end{center}
{\small \textbf{Figure 22.} \textbf{How to read the figure.} At each displayed mass, the black signal-MC distribution is excluded from model construction and then used to check the prediction. Blue interpolates the two adjacent samples, including core and tails. Red uses the equal-weight common shape from all other anchors. Both models predict the target center and width from its neighbors. Upper panels show the core, and lower panels show the tails. All densities retain the full selected-event normalization; no fit-window renormalization or uncertainty bands are applied.}\par\medskip

The signal-MC sample at the tested mass is removed from both template construction and the prediction of its core parameters. It is then used as the independent shape target. For example, the 100 MeV prediction uses 80 and 120 MeV; the 160 MeV prediction uses 140 and 180 MeV. This is more demanding than reproducing an anchor, where interpolation equals the input histogram by construction.

Neighboring-template interpolation captures the mass dependence of the tails more closely than the common shape in this diagnostic. The largest residual occurs near the lower end of the tested range: the 100 MeV prediction must bridge the rapidly changing 80--120 MeV distributions. Similar-looking cores therefore do not justify assuming exact tail agreement.

These comparisons use the full selected normalization. Neither a good core fit nor a smaller shape discrepancy by itself establishes improved sensitivity; the independent toy study below tests that next.
\clearpage\section*{Appendix E. Quantify interpolation errors across the range}
\begin{center}\includegraphics[width=\linewidth,height=4.5in,keepaspectratio]{../figures/../appendix_v637/figures/v637_shape_closure_metrics.pdf}\end{center}
{\small \textbf{Figure 23.} \textbf{How to read the figure.} Each target mass is removed before predicting its center, width and distribution. The upper panels measure core-parameter errors. The lower-left panel gives the maximum absolute CDF difference, which compares integrated probability throughout the stored mass range. The lower-right panel compares probability in the same [-4,3] window defined using the excluded MC core only for evaluation. Smaller absolute differences indicate better agreement. These are deterministic shape checks; finite signal-MC and fitted-core uncertainties are not propagated, so no bars are shown.}\par\medskip

For the seven omitted masses from 100 to 220 MeV, the largest absolute cumulative-probability difference is 0.0040--0.0276 for neighboring-template interpolation, compared with 0.0260--0.0421 for the common shape. This metric asks how far the predicted integrated probability can depart from the simulated distribution at any stored mass threshold. It is not a goodness-of-fit probability or a confidence interval.

The comparison at 100 MeV has a difference of about 0.0276: an integrated probability can be wrong by about 2.8 percentage points. This remaining modeling difference is relevant even though the core looks close. Finite signal-MC statistics and the core-fit uncertainties are not propagated into the new toy calibration; the template inputs are treated as fixed.

The package checks nonnegative bin probabilities, monotone cumulative distributions, unit total probability including outside-support categories, and exact reproduction at the retained anchors. Interpolation is allowed only inside 80--240 MeV. A curve drawn at an intermediate mass is a model prediction, not new detector simulation.
\clearpage\section*{Appendix F. Choose the window before testing it}

Three simulated masses, 100, 160 and 220 MeV, span the interior of the interpolation range. Each is removed from the common and neighboring template inputs. The direct signal-MC extraction is a best-case benchmark at a known mass; it is not eligible to win an interpolation procedure that must also work between samples.

The original shifted Gaussian uses its established center and reference width, with both fit and GP exclusion at $\pm2.25\sigma_{\rm ref}$. All new windows use the center and core width predicted from neighboring samples. The Gaussian in the starter window retains its reference Gaussian width. Common, neighboring and direct MC templates in the same starter window isolate the shape comparison.
\begin{center}\small\begin{tabular}{lll}\toprule
Extraction template & Fit interval $F$ & GP exclusion $G$\\\midrule
Original shifted Gaussian & $\pm2.25\sigma_{\rm ref}$ & $\pm2.25\sigma_{\rm ref}$\\
Gaussian / common / neighbors / direct MC & $[-4,3]$ & $[-4,3]$\\
Neighbors: shortened & $[-2.5,2.25]$ & $[-2.5,2.25]$\\
Neighbors: intermediate & $[-3,2.5]$ & $[-3,2.5]$\\
Neighbors: one-unit extension & $[-5,4]$ & $[-5,4]$\\
Neighbors: two-unit extension & $[-6,5]$ & $[-6,5]$\\
Neighbors: modest unused buffer & $[-3,2.5]$ & $[-4,3]$\\
Neighbors: wide unused buffer & $[-4,3]$ & $[-6,5]$\\
\bottomrule\end{tabular}\end{center}
Except for the first row, intervals are in $u$ and the actual bin boundaries are saved in the geometry table. This gives eleven candidates. The training data and the signal-fit data are disjoint for every candidate.

\textbf{Tuning.} Use 100 GP mean-source experiments with background only and with full signal MC at $A=3s_0$. The scale $s_0$ is the archived mean Gaussian yield error from an independent 100-toy background-only pilot at that mass. The expected injected yield is the same for every candidate. Poisson signal increments and background draws are paired across methods. Measure
\[
 R=\frac{\operatorname{mean}(\widehat A_{\rm signal+background}-\widehat A_{\rm background})}{A},
 \qquad W=\frac{\operatorname{SD}(\widehat A_{\rm background})}{R}.
\]
A smaller $W$ means less fitted-yield noise after correcting the mean response to added signal. Choose one candidate for all three masses by minimizing their geometric mean $W$. This is a sensitivity proxy, not an upper limit. The choice is frozen before calibration and evaluation. The common shape with $[-4,3]$ wins; its tuning ratio to the original Gaussian is 0.9917, only a 0.8\% reduction.

\textbf{Independent test.} For each background source, generate 100 new calibration experiments and 100 separate evaluation experiments. Calibration spans $A/s_0=0,0.5,1,2,3,4,5,6,8,10,12,16$; evaluation uses $0,3,5$. The selected common shape, original Gaussian, and three other starter-window shapes are retained regardless of their ranking. The functional form background source supplies a separate stress test with its own calibration; it is not a GP-calibrated test under a mismatched source.
\clearpage\section*{Appendix F. Shorter and wider windows: the measured tradeoff}
\begin{center}\includegraphics[width=\linewidth,height=5.0in,keepaspectratio]{../figures/../appendix_v637/figures/v637_window_scan.pdf}\end{center}
{\small \textbf{Figure 24.} \textbf{How to read the figure.} Every curve uses full v16 TC signal-MC injections and an extraction template predicted from neighboring masses. Upper panel: fraction of all selected signal events entering GP training bins. Lower panel: background-only fitted-yield scatter divided by the signal response, relative to the original Gaussian result at that mass. Bars are 95\% percentile intervals from 2,000 whole-toy bootstrap resamples of the 100 tuning experiments, preserving mass and method correlations. $s_0$ is the archived Gaussian yield-error scale. The last two choices leave an unused buffer between fit and training. The upper panel has no uncertainty bars because the signal histogram is treated as fixed.}\par\medskip

The proposed $[-4,3]$ interval leaves 20.0\%, 16.2\% and 19.5\% of the full selected signal in GP training at 100, 160 and 220 MeV. Extending the fit and exclusion to $[-6,5]$ reduces those fractions to 14.6\%, 10.9\% and 12.4\%. They are still not negligible. All these tails are present in the simulated data.

The two-unit extension does not improve the tuning precision: its three-mass ratio is 1.0109. Excluding that wide interval while fitting only $[-4,3]$ is more costly, with ratio 1.1846. In these tests, removing background information without fitting the additional bins is worse than using the wider fit. No tested interval achieves both negligible leakage and a demonstrated sensitivity gain.
\clearpage\section*{Appendix G. Better yield response, modest sensitivity changes}
\begin{center}\includegraphics[width=\linewidth,height=4.2in,keepaspectratio]{../figures/../appendix_v637/figures/v637_independent_response.pdf}\end{center}
{\small \textbf{Figure 25.} \textbf{How to read the figure.} Each point uses 100 independent evaluation experiments with expected full selected yield $A=3s_0$. Left: mean added fitted yield divided by $A$; bars are sample standard errors of the paired signal-induced change. Right: background-only yield noise divided by the response, relative to the original shifted Gaussian; bars are 95\% intervals from 2,000 whole-toy bootstrap resamples. Resampling preserves mass, strength and method correlations; sources use separate streams. The pilot scale is fixed. The common-shape window was selected on a separate tuning sample. A response closer to one does not by itself mean a smaller upper limit.}\par\medskip
\begin{center}\small\begin{tabular}{lrrr}\toprule
Mass [MeV] & Common response $R$ & Noise ratio & 95\% bootstrap interval\\\midrule
100 & 1.003 & 0.969 & [0.937, 1.000]\\
160 & 1.027 & 1.013 & [0.979, 1.048]\\
220 & 0.857 & 0.952 & [0.915, 0.992]\\
\bottomrule\end{tabular}\end{center}
For the selected common template, the geometric-mean noise ratio across the three masses is 0.9774 [0.9579, 0.9977] on the GP mean source and 0.9696 [0.9451, 0.9919] on the functional form source. These are conditional reductions of about 2--3\% in this precision diagnostic. The mass-by-mass result is mixed, including a slight increase at 160 MeV on the GP source.

The direct MC template still returns responses of about 0.959, 0.950 and 0.909 on the GP source. Even the exact signal shape in the fitted bins does not eliminate absorption of signal in the sidebands. The common shape's response can also exceed one because its mass-dependent tail fraction is imperfect. Calibration remains necessary.
\clearpage\section*{Appendix G. Toy-calibrated limits and signal detection}
\begin{center}\includegraphics[width=\linewidth,height=4.15in,keepaspectratio]{../figures/../appendix_v637/figures/v637_limits_and_detection.pdf}\end{center}
{\small \textbf{Figure 26.} \textbf{How to read the figure.} Calibration and evaluation each contain 100 experiments per source. Left: median largest accepted $A/s_0$ from the discrete yield grid, calculated among nonempty accepted sets; no error bars are assigned to these coarse medians. Empty counts are listed below. Right: fraction rejecting background only when full signal MC at $A=3s_0$ is injected. Bars are exact two-sided 95\% Clopper--Pearson intervals with denominator 100, conditional on the saved calibration sample. Calibration uses the same background source and direct signal-MC distribution as evaluation. Curves at common coordinates can overlap. These are fixed-mass local tests; a 10\% rejection threshold is not a discovery criterion.}\par\medskip
\begin{center}\small\begin{tabular}{lrrrrr}\toprule
Source & MeV & Gaussian $U/s_0$ & Empty / 100 & Common $U/s_0$ & Empty / 100\\\midrule
GP mean & 100 & 1.50 & 8 & 2.00 & 8\\
GP mean & 160 & 1.00 & 6 & 1.00 & 7\\
GP mean & 220 & 2.00 & 8 & 2.00 & 2\\
Functional & 100 & 1.00 & 11 & 1.00 & 11\\
Functional & 160 & 1.00 & 6 & 1.00 & 5\\
Functional & 220 & 2.00 & 5 & 1.00 & 5\\
\bottomrule\end{tabular}\end{center}
The endpoint comparison does not show uniformly improved upper limits. At 100 MeV on the GP source the selected common shape has a larger median endpoint, while several other comparisons tie on this coarse grid. Calibration sampling and grid spacing limit what these medians can resolve. The 2--3\% change in the precision diagnostic must not be restated as a measured 2--3\% limit improvement.
\clearpage\section*{Appendix G. What is calibrated, and what is not?}

\textbf{Upper-limit construction.} For each tested expected yield $A$, fit 100 independent calibration spectra and compare the evaluation fitted yield with that distribution:
\[
 p_A=\frac{1+\#\{\widehat A_{{\rm cal},A}\le\widehat A_{\rm eval}\}}{101}.
\]
Retain every grid value with $p_A>0.10$. The stored result is the complete accepted set. Its largest member is a grid endpoint, not an interpolated continuous upper limit. An empty set means none of the tested yields was accepted; it is never reported as a physical zero limit. A hole or acceptance of the largest tested yield would also need separate treatment.

Across all 9,000 primary evaluation fits, 205 accepted sets are empty, none has holes, and none reaches the top of the grid. The preceding table gives the relevant background-only empty counts beside each displayed endpoint median. All 100 experiments remain in rejection-rate and detection-rate denominators; only the endpoint median conditions on a nonempty set.

\textbf{Background-only test.} To test whether a fitted yield is unusually high under background only, reverse the rank comparison:
\[
 p_0=\frac{1+\#\{\widehat A_{{\rm cal},0}\ge\widehat A_{\rm eval}\}}{101}.
\]
Reject background only at $p_0\le0.10$. With exchangeable continuous calibration and evaluation statistics, this rank rule has marginal rejection probability $10/101$ under the tested hypothesis; ties make it conservative. Each realized calibration sample gives a fluctuating threshold. Binomial intervals for the evaluation counts condition on that saved threshold and do not include drawing a replacement calibration sample.

For the selected common shape, the background-only rejection counts at 100, 160 and 220 MeV are 12, 5, 9 out of 100 on the GP mean source. The corresponding functional-source counts are 4, 10, 9. These finite-sample checks are distinct from signal detection.

\textbf{Why matching the calibration truth matters.} The primary comparison calibrates every extraction method using the full direct MC distribution at the tested mass. That makes the comparison fair at those known samples, but assumes the true signal distribution is available for calibration. At an unsimulated mass, calibration must instead rely on an interpolated shape. The next page deliberately changes that assumption and checks the resulting rejection rate against the removed signal-MC sample.

\textbf{Uncertainties.} Mean bars are sample standard errors. Widths and precision ratios use 2,000 whole-toy bootstrap resamples that preserve masses, strengths and policy correlations; sources use separate streams. Binomial counts retain denominator 100 and exact two-sided 95\% Clopper--Pearson intervals. Finite signal-MC and fitted-core uncertainties remain fixed in this bounded study.
\clearpage\section*{Appendix G. Calibrate with interpolation, test against signal MC}
\begin{center}\includegraphics[width=\linewidth,height=3.6in,keepaspectratio]{../figures/../appendix_v637/figures/v637_calibration_mismatch.pdf}\end{center}
{\small \textbf{Figure 27.} \textbf{How to read the figure.} Calibration now injects the omitted-mass common or neighboring prediction, while evaluation still injects the full direct signal-MC sample. Each method uses 100 fresh calibration experiments on the GP mean source and the same 100 independent evaluation experiments used above. Points show how often the true injected grid yield is rejected at $A=3s_0$ and $5s_0$. Bars are exact two-sided 95\% Clopper--Pearson intervals, denominator 100, conditional on that calibration sample. The dashed $10/101$ reference applies under exchangeability; shape mismatch removes that guarantee. The calibration stream differs from the primary comparison, so changes also contain finite-calibration fluctuations.}\par\medskip

At 100, 160 and 220 MeV, calibrating with the common shape rejects the true injected yield in 6, 8 and 12 of 100 experiments at $A=3s_0$, and in 5, 4 and 13 at $A=5s_0$. Neighboring-template calibration gives 6, 13 and 5 at $3s_0$, and 6, 10 and 4 at $5s_0$. Every displayed 95\% interval contains $10/101$. There are no empty sets at these two injected strengths, no holes, and no top-grid censoring.

This is a useful closure check for the practical case in which the local direct MC template is unavailable. It does not prove exact coverage or resolve small differences between the two interpolation prescriptions. Only three interior masses, two injected strengths, one background source and the fixed selected signal-MC inputs were tested in this additional diagnostic. The 90 MeV curve remains a prediction without an independent detector-MC target.

The additional study has its own frozen protocol and random-number namespace. It does not change the tuning selection or reuse the evaluation experiments to choose another window.
\clearpage\section*{Appendix H. Conclusions, sources and reproducibility}

\textbf{What v6.3.7 establishes.} The v16 neighboring-template construction predicts the full shape more accurately than a universal aligned shape in the omitted-mass checks. Both models predict the core center and width from the same neighbors. MC-informed extraction substantially improves the relation between added full signal yield and fitted yield. The selected starter-window common shape gives a modest average precision gain in independent toys, but the tested upper limits and detection fractions do not establish a uniform sensitivity improvement across masses.

\textbf{What the window study changes.} The proposed $[-4,3]$ region is a useful tested starting point, not a negligible-leakage prescription. Even the $[-6,5]$ extension leaves roughly 11--15\% of selected signal in training bins at these masses. Broadening only the exclusion while keeping a narrow fit can substantially worsen precision. The original observed-data analysis is therefore retained; the new templates and comparisons are supplied as reproducible candidates for a full mass-grid study, including selection equivalence and template uncertainty.

\textbf{Numerical record.} All 58,800 toy extraction fits passed the saved numerical checks: 6,600 tuning, 36,000 primary calibration, 9,000 independent evaluation and 7,200 interpolation-calibration fits. Thirty additional deterministic profile-limit calculations passed and are saved as asymptotic model diagnostics; they are not the toy-rank results shown here. The statistics review independently recomputed the rank tests, binomial intervals, cohort separation, pairing, geometry and normalization. No new observed-data fit or UC upper limit was made.

\textbf{Reproduce the study.} The package includes the v16 histograms, fitted core summaries, archived background spectra and kernel parameters, complete toy-fit rows, signal-category draws, random seeds, frozen candidate selection and accepted yield sets. Input hashes identify the analyzed data. The scripts \path{validate_shapes.py}, \path{run_study.py}, \path{calibration_mismatch.py}, \path{summarize_global.py}, \path{audit_statistics.py}, \path{make_figures.py} and \path{build_report.py} rebuild the numerical checks, figures and report locally. Checkpoints verify their protocol and row hashes before reuse. The updated v6.3.6 package embeds this study in \path{appendix_v637/}; previous releases are preserved separately.

\textbf{Literature and its role.}

[1] M. Baak, S. Gadatsch, R. Harrington and W. Verkerke, \emph{Interpolation between multi-dimensional histograms using a new non-linear moment morphing method}, NIM A771 (2015) 39--48, \href{https://arxiv.org/abs/1410.7388}{arXiv:1410.7388}. This motivates aligning location and width before mixing templates and testing predictions against omitted simulations. Our empirical core-aligned construction is related, but is not an exact implementation of moment morphing.

[2] M. Frate et al., \emph{Modeling Smooth Backgrounds and Generic Localized Signals with Gaussian Processes}, \href{https://arxiv.org/abs/1709.05681}{arXiv:1709.05681}. Its discussion of regularized GP likelihoods and injected-signal power motivates ensemble calibration of the background-plus-signal procedure; it does not validate the present kernel choices.

[3] G. Punzi, \emph{Sensitivity of searches for new signals and its optimization}, \href{https://arxiv.org/abs/physics/0308063}{arXiv:physics/0308063}. Search optimization should address detection and exclusion. Here independent toys evaluate those quantities after tuning; no counting-experiment approximation replaces the correlated GP calculation.

[4] G. Cowan et al., \emph{Asymptotic formulae for likelihood-based tests of new physics}, EPJ C71 (2011) 1554, \href{https://arxiv.org/abs/1007.1727}{arXiv:1007.1727}. This supplies the context for the separately saved asymptotic profile-limit diagnostics, not a coverage guarantee for these two-stage GP fits.
\clearpage\section*{Appendix I. How wide is the proposed blind window?}

This follow-up extends the comparison to nine v16 TC signal-MC masses from 80 to 240 MeV. The requested interval is $u\in[-4,3]$, where $u=(m_{\rm rec}-c)/\sigma_{\rm core}$. Its width is $7\sigma_{\rm core}$, while the historical interval has width $4.5\sigma_m$. Here $\sigma_m=\sigma_{\rm ref}$ is the historical analysis resolution; it is not the fitted signal-MC core width.
\begin{center}\includegraphics[width=\linewidth,height=3.15in,keepaspectratio]{../figures/../appendix_leakage/note_figures/leakage_window_width.pdf}\end{center}
{\small \textbf{Figure 28.} \textbf{How to read the figure.} Left: the requested lower and upper offsets expressed in historical $\sigma_m$ units; dashed lines show the historical offsets $\pm2.25$. Offsets refer to each window's own center. Right: percentage increase in total continuous width and in the actual excluded-bin width. Interior core parameters are predicted from neighboring signal-MC samples with the tested sample omitted, as in v6.3.7; the 80 and 240 MeV endpoints use their directly fitted cores. These fixed-template geometry calculations have no uncertainty bars. The analysis spectrum is binned, so bin counts need not increase by the continuous-width percentage.}\par\medskip
\begin{center}\small\begin{tabular}{lrrrr}\toprule
Mass [MeV] & Lower / $\sigma_m$ & Upper / $\sigma_m$ & Wider by & Old / new bins\\\midrule
80 & -2.86 & +2.15 & 11.3\% & 17 / 18\\
100 & -2.92 & +2.19 & 13.5\% & 19 / 21\\
120 & -2.91 & +2.18 & 13.1\% & 21 / 23\\
140 & -2.89 & +2.17 & 12.4\% & 25 / 27\\
160 & -2.88 & +2.16 & 12.0\% & 28 / 31\\
180 & -2.87 & +2.15 & 11.6\% & 31 / 35\\
200 & -2.94 & +2.20 & 14.2\% & 36 / 41\\
220 & -3.00 & +2.25 & 16.7\% & 41 / 48\\
240 & -3.08 & +2.31 & 19.8\% & 47 / 55\\
\bottomrule\end{tabular}\end{center}
The requested interval is 11.3--19.8\% wider continuously across this range, predominantly on the low-mass side. At the three original toy-test masses, 100, 160 and 220 MeV, it is 13.5\%, 12.0\% and 16.7\% wider. A seven-unit interval in $u$ is therefore much less than 56\% wider than the historical interval: the units differ because $\sigma_{\rm core}<\sigma_m$.
\clearpage\section*{Appendix I. How much signal enters GP training?}
\begin{center}\includegraphics[width=\linewidth,height=3.4in,keepaspectratio]{../figures/../appendix_leakage/note_figures/leakage_probability.pdf}\end{center}
{\small \textbf{Figure 29.} \textbf{How to read the figure.} Left: full selected signal probability in GP training bins for the historical and starter windows. Solid curves use actual v16 TC signal MC; dashed curves use a hypothetical shifted Gaussian with the historical analysis width $\sigma_m$. Right: MC-to-Gaussian leakage ratio within the same window and recorded spectrum. This compares signal distributions, not the ability of a Gaussian fit to remove actual MC tails. All probabilities retain full selected normalization. Signal outside the background spectrum is stored separately and is not counted as GP training contamination. No uncertainty bars are shown; finite signal-MC and core-fit uncertainties are fixed in this diagnostic.}\par\medskip
\begin{center}\small\begin{tabular}{lrrrr}\toprule
Mass [MeV] & MC old [\%] & MC new [\%] & Shifted Gaussian, new [\%] & Reduction [pp]\\\midrule
80 & 29.02 & 27.90 & 2.19 & 1.12\\
100 & 21.77 & 19.99 & 1.39 & 1.77\\
120 & 19.61 & 17.61 & 1.91 & 2.01\\
140 & 17.76 & 16.16 & 1.74 & 1.60\\
160 & 18.22 & 16.19 & 1.86 & 2.03\\
180 & 19.51 & 17.06 & 1.66 & 2.45\\
200 & 20.59 & 17.99 & 1.51 & 2.61\\
220 & 22.65 & 19.52 & 1.49 & 3.13\\
240 & 24.96 & 21.66 & 1.29 & 3.30\\
\bottomrule\end{tabular}\end{center}
The wider starter window removes an additional 1.12--3.30 percentage points (pp) of the full selected MC signal from GP training. Nevertheless, 16.16--27.90\% remains, compared with 1.29--2.19\% for the shifted Gaussian in the same bins: approximately 8.7--16.8 times more leakage. For reference, an ideal continuous Gaussian has 2.445\% outside $\pm2.25\sigma_m$; actual binned masks differ slightly.

\textbf{Leakage is not the fitted-yield loss.} A broad tail can enter training without being absorbed one-for-one into a narrow signal estimate. To measure the effect on the fit, keep its input counts fixed and change only the contamination of the GP training data. The next two pages implement that controlled comparison.
\clearpage\section*{Appendix I. Fit impact on the GP mean background source}

Reuse the saved full signal-MC injections at expected yield $A=3s_0$, where $s_0$ is the archived Gaussian background-only yield-error scale. In the control, subtract the known injected signal only from GP training bins. Keep the signal-fit counts, extraction template, window and kernel settings identical, then recompute the GP mean and covariance. This isolates the effect of contaminated training data.
\begin{center}\includegraphics[width=\linewidth,height=4.3in,keepaspectratio]{../figures/../appendix_leakage/note_figures/leakage_fit_gp_mean.pdf}\end{center}
{\small \textbf{Figure 30.} \textbf{How to read the figure.} Results use the GP mean background source. Lines come from deterministic mean-count spectra at nine masses; markers come from 100 paired saved evaluation experiments at 100, 160 and 220 MeV. Upper panel: response $R=\operatorname{mean}(\widehat A_{s+b}-\widehat A_b)/A$, with contaminated and clean training. Lower left: response loss $\operatorname{mean}(\widehat A_{\rm clean}-\widehat A_{\rm contaminated})/A$. Mean bars are sample standard errors of paired responses or differences. Lower right: extra response-scaled background-only noise, $100(R_{\rm clean}/R_{\rm contaminated}-1)$, with 95\% percentile intervals from 2,000 whole-toy bootstrap resamples preserving masses and methods; sources use separate streams. Background-only fits are unchanged by the control. Blue uses direct MC extraction with the starter window; orange uses the shifted Gaussian with the historical window. Small bars can be hidden by markers. These precision changes are not upper-limit improvements. The table below uses the direct MC template with the starter window.}\par\medskip
\begin{center}\small\begin{tabular}{lrrrr}\toprule
MeV & Contaminated $R$ & Clean $R$ & Yield lost [\% of $A$] & Extra noise [\%; 95\% interval]\\\midrule
100 & 0.959 & 1.001 & $4.19\pm0.02$ & 4.37 [4.32, 4.41]\\
160 & 0.950 & 0.997 & $4.71\pm0.03$ & 4.96 [4.90, 5.01]\\
220 & 0.909 & 1.000 & $9.13\pm0.05$ & 10.05 [9.92, 10.18]\\
\bottomrule\end{tabular}\end{center}Removing contamination restores the direct-MC response to approximately one; the measured loss is 4.2--9.1\% of the full injected yield. The shifted Gaussian retains a response deficit because its shape and yield convention do not match the MC distribution.\clearpage\section*{Appendix I. Fit impact on the functional form background source}

Reuse the saved full signal-MC injections at expected yield $A=3s_0$, where $s_0$ is the archived Gaussian background-only yield-error scale. In the control, subtract the known injected signal only from GP training bins. Keep the signal-fit counts, extraction template, window and kernel settings identical, then recompute the GP mean and covariance. This isolates the effect of contaminated training data.
\begin{center}\includegraphics[width=\linewidth,height=4.3in,keepaspectratio]{../figures/../appendix_leakage/note_figures/leakage_fit_functional.pdf}\end{center}
{\small \textbf{Figure 31.} \textbf{How to read the figure.} Results use the functional form background source. Lines come from deterministic mean-count spectra at nine masses; markers come from 100 paired saved evaluation experiments at 100, 160 and 220 MeV. Upper panel: response $R=\operatorname{mean}(\widehat A_{s+b}-\widehat A_b)/A$, with contaminated and clean training. Lower left: response loss $\operatorname{mean}(\widehat A_{\rm clean}-\widehat A_{\rm contaminated})/A$. Mean bars are sample standard errors of paired responses or differences. Lower right: extra response-scaled background-only noise, $100(R_{\rm clean}/R_{\rm contaminated}-1)$, with 95\% percentile intervals from 2,000 whole-toy bootstrap resamples preserving masses and methods; sources use separate streams. Background-only fits are unchanged by the control. Blue uses direct MC extraction with the starter window; orange uses the shifted Gaussian with the historical window. Small bars can be hidden by markers. These precision changes are not upper-limit improvements. The table below uses the direct MC template with the starter window.}\par\medskip
\begin{center}\small\begin{tabular}{lrrrr}\toprule
MeV & Contaminated $R$ & Clean $R$ & Yield lost [\% of $A$] & Extra noise [\%; 95\% interval]\\\midrule
100 & 0.957 & 0.999 & $4.20\pm0.02$ & 4.39 [4.35, 4.43]\\
160 & 0.953 & 1.000 & $4.76\pm0.02$ & 4.99 [4.95, 5.04]\\
220 & 0.908 & 1.000 & $9.25\pm0.06$ & 10.19 [10.04, 10.33]\\
\bottomrule\end{tabular}\end{center}The functional form source gives similar results. Removing known injected signal is a simulation diagnostic, not a prescription for data or a demonstrated upper-limit gain. All 3,000 control fits, 270 deterministic fits and 30 replays passed. No new toy draws or upper limits were produced. Source and validation are in \path{appendix_leakage/}; earlier results remain unchanged.\clearpage\section*{Appendix I. Does the Gaussian center change the comparison?}

Yes: the appropriate Gaussian reference follows the displaced signal core. The preceding leakage comparison already uses this shifted Gaussian. This added check holds the historical analysis width $\sigma_m$ and every blind/training bin fixed, and changes only the Gaussian center. For generated mass $M$, the two centers are
\[
 \mu_{\rm central}=M,\qquad
 \mu_{\rm shifted}=M-3.22433\,{\rm MeV}-2.21399\,{\rm MeV}\ln\!\left(\frac{M}{150\,{\rm MeV}}\right).
\]
\begin{center}\includegraphics[width=\linewidth,height=3.2in,keepaspectratio]{../figures/../appendix_leakage/note_figures/gaussian_center_leakage_comparison.pdf}\end{center}
{\small \textbf{Figure 32.} \textbf{How to read the figure.} Full selected signal probability entering the recorded GP training bins at nine v16 TC masses. Orange: Gaussian centered at the generated mass. Blue: shifted Gaussian, the reference already used in Figure 29. Both have width $\sigma_m$. Gray: unchanged signal MC. Left uses the historical blind window centered at $\mu_{\rm shifted}$; right uses the starter window around the predicted MC core. Window positions stay fixed between curves, so the difference between Gaussian curves isolates their center. Probabilities outside the recorded spectrum are retained separately. No uncertainty bars are shown: these are fixed-template probability integrals, with no new fits or toys.}\par\medskip
\begin{center}\small\begin{tabular}{lrrrr}\toprule
Mass [MeV] & Old: central [\%] & Old: shifted [\%] & New: central [\%] & New: shifted [\%]\\\midrule
80 & 6.40 & 2.02 & 10.46 & 2.19\\
100 & 8.71 & 2.10 & 8.67 & 1.39\\
120 & 7.99 & 2.53 & 11.54 & 1.91\\
140 & 7.80 & 1.93 & 10.84 & 1.74\\
160 & 8.08 & 2.22 & 10.73 & 1.86\\
180 & 8.67 & 2.71 & 8.57 & 1.66\\
200 & 7.50 & 2.47 & 7.40 & 1.51\\
220 & 6.94 & 2.48 & 6.81 & 1.49\\
240 & 5.59 & 2.31 & 5.44 & 1.29\\
\bottomrule\end{tabular}\end{center}
At 100, 160 and 220 MeV, the starter-window leakage is 8.67\%, 10.73\% and 6.81\% for the central-mass Gaussian, versus 1.39\%, 1.86\% and 1.49\% for the shifted Gaussian. The MC leakage remains 19.99\%, 16.19\% and 19.52\%. The shifted reference makes the excess MC-tail contamination more apparent.

Existing fit-impact controls already use shifted Gaussian extraction and are unchanged. Moving a Gaussian together with its continuous symmetric window would instead keep leakage outside $\pm2.25\sigma_m$ fixed at 2.445\%.

\end{document}
